\subsection{Extraneous Solutions to a Non-Linear System} 
In the previous example, we solved the system
\[\begin{systemL}
x^2+xy+y^2=19&\\
y^2-x^2=5&
\end{systemL}\]
and got \makeblanksmall (apparent) solutions:
\fakeit{%
\begin{math}
\begin{array}{r@{{}={}}lr@{{}={}}lr@{{}={}}lr@{{}={}}l}
(x,y)&\left(2,3\right)&
(x,y)&\left(2,-3\right)&
(x,y)&\left(-2,3\right)&
(x,y)&\left(-2,-3\right)\\
(x,y)&\left(\tfrac{7}{\sqrt{3}},\tfrac{8}{\sqrt{3}}\right)&
(x,y)&\left(\tfrac{7}{\sqrt{3}},-\tfrac{8}{\sqrt{3}}\right)&
(x,y)&\left(-\tfrac{7}{\sqrt{3}},\tfrac{8}{\sqrt{3}}\right)&
(x,y)&\left(-\tfrac{7}{\sqrt{3}},-\tfrac{8}{\sqrt{3}}\right)
\end{array}
\end{math}
}
But when we graph this system, we see:
\begin{icpic}[x=.5\linespace,y=.5\linespace]
\setgraphsquare{6}
\GraphLarge
\draw[graph,domain={-sqrt(31)}:{sqrt(31)},samples=100] plot (\x,{sqrt(5+\x*\x)});
\draw[graph,domain={-sqrt(31)}:{sqrt(31)},samples=100] plot (\x,{-sqrt(5+\x*\x)});
\draw[graphend,domain=0:360,samples=360] plot 
          ({sqrt(19/3)*cos(\x)+sqrt(19)*sin(\x)},{sqrt(19/3)*cos(\x)-sqrt(19)*sin(\x)});
\end{icpic}
These curves only cross \makeblanksmall times! We must have \makeblank[1in] solutions. Steps that introduce extraneous solutions include:
\begin{itemize}
\item Multiplying both sides of an equation by a \makeblank[1in] expression% (might multiply by \makeblanksmall)
\item Taking an \makeblank[1in] power of an \makeblank[1in] radical% (might lose \makeblank[1in] information)
\item Combining \makeblank[1in] expressions% (might affect the \makeblank[1in])
\end{itemize}
In our example, we \makeblank. So we must check for extraneous solutions. (When in doubt, check!)
\begin{example}[continued]
Check for extraneous solutions to
\[\begin{systemL}
x^2+xy+y^2=19&\\
y^2-x^2=5&
\end{systemL}\]
\begin{note}
We used $(II)$ to find the possible solutions, using steps that did \emph{not} risk introducing extraneous solutions. We need to check which solutions also work in $(I)$.
\end{note}
\fakeit{%
\begin{math}\begin{array}{ror@{}rororror@{}roror}
( 2)^2&+&( 2)&( 3)&+&( 3)^2&=&\makeblanksmall&
\left(\tfrac{7}{\sqrt{3}}\right)^2&
+&\left(\tfrac{7}{\sqrt{3}}\right)&\left(\tfrac{8}{\sqrt{3}}\right)&
+&\left(\tfrac{8}{\sqrt{3}}\right)^2&
=&\makeblanksmall\\
( 2)^2&+&( 2)&(-3)&+&(-3)^2&=&\makeblanksmall&
\left(\tfrac{7}{\sqrt{3}}\right)^2&
+&\left(\tfrac{7}{\sqrt{3}}\right)&\left(-\tfrac{8}{\sqrt{3}}\right)&
+&\left(-\tfrac{8}{\sqrt{3}}\right)^2&
=&\makeblanksmall\\
(-2)^2&+&(-2)&( 3)&+&( 3)^2&=&\makeblanksmall&
\left(-\tfrac{7}{\sqrt{3}}\right)^2&
+&\left(-\tfrac{7}{\sqrt{3}}\right)&\left(\tfrac{8}{\sqrt{3}}\right)&
+&\left(\tfrac{8}{\sqrt{3}}\right)^2&
=&\makeblanksmall\\
(-2)^2&+&(-2)&(-3)&+&(-3)^2&=&\makeblanksmall&
\left(-\tfrac{7}{\sqrt{3}}\right)^2&
+&\left(-\tfrac{7}{\sqrt{3}}\right)&\left(-\tfrac{8}{\sqrt{3}}\right)&
+&\left(-\tfrac{8}{\sqrt{3}}\right)^2&
=&\makeblanksmall
\end{array}\end{math}}
\end{example}